% =========================
% Explanation: Benchmarks / Problem Instances (informal, readable)
% =========================
\subsection*{Explanation of Benchmarks / Problem Instances (informal)}
This section specifies what problems we run the optimizers on and what exactly an objective evaluation means.
In benchmarking, these definitions are critical: they determine what is comparable and what claims are valid.

\begin{itemize}
    \item \textbf{What is a ``problem instance''?}
    A benchmark \emph{problem instance} is a fully specified optimization task given to an HPO/NAS method.
    In practice, an instance can be written as
    \[
        \text{instance} =
        (\text{dataset},\ \text{split},\ \text{preprocessing},\ \text{training recipe},\ \Lambda,\ \text{objective definition}),
    \]
    where \(\Lambda\) is the search space and the objective definition specifies how a configuration \(\lambda\) is mapped to objective value(s).

    \item \textbf{Datasets and splits.}
    For each dataset we specify:
    (i) dataset name and version,
    (ii) the train/validation/test split policy (or cross-validation folds),
    and (iii) any access constraints (e.g., an official test server).
    Fixing splits is important because different splits can change the objective landscape and confound algorithm comparisons.
    We therefore either use canonical public splits or publish a split manifest so the same examples are used across methods.

    \item \textbf{Preprocessing and evaluation pipeline.}
    We specify the preprocessing steps applied to the raw data (e.g., normalization, augmentation, tokenization, feature extraction)
    and define the full evaluation pipeline that turns \(\lambda\) into objective value(s).
    The pipeline typically includes model instantiation, training under a fixed recipe, validation scoring,
    and (optionally) measurement of secondary objectives such as FLOPs, latency, memory, energy, or parameter count.
    For fairness, this pipeline is identical for all optimizers.

    \item \textbf{Objective function(s) \(f(\lambda)\) and \(\mathbf{F}(\lambda)\).}
    The objective function formalizes what is being optimized.
    In the single-objective case, we define a scalar objective
    \[
        f(\lambda) \in \mathbb{R},
    \]
    such as validation error or negative log-likelihood.
    In the multiobjective case, we define a vector-valued objective
    \[
        \mathbf{F}(\lambda) = (f_1(\lambda),\ldots,f_m(\lambda)) \in \mathbb{R}^m,
    \]
    where \(m\) is the number of objectives.
    We also specify the optimization direction (minimize vs.\ maximize) and any transformations used to convert maximization metrics to minimization.

    \item \textbf{Normalization/scaling for Pareto metrics (multiobjective).}
    Pareto-quality metrics such as hypervolume and the \(\epsilon\)-indicator can be sensitive to objective scaling.
    Therefore, when we compute these metrics, we specify a fixed monotone normalization rule
    (e.g., affine scaling to \([0,1]\) using predefined bounds) and apply it identically across methods on each instance.
    This ensures metric differences reflect Pareto quality rather than arbitrary units.

    \item \textbf{Search space \(\Lambda\) and constraints.}
    The benchmark definition includes the complete search space:
    all hyperparameters/architecture choices, their types (continuous/integer/categorical/structured),
    ranges, and any conditional dependencies.
    Constraints define which configurations are valid; invalid configurations must be handled consistently
    (repair, rejection, or penalty).
    Sharing the same \(\Lambda\) is a core fairness requirement: no method may use a different space.

    \item \textbf{Representativeness and claim support.}
    The benchmark suite should reflect the intended use case of the proposed method.
    We justify why the chosen tasks/datasets (including their sizes, modalities, and regimes) are representative of the paper’s claims.
    Importantly, conclusions are bounded to the tested instances: performance on this suite supports claims about this regime,
    but does not automatically imply performance on substantially different tasks unless those are included explicitly.
\end{itemize}
